\subsection{一个例子}\label{example}

本节详细算一个扭余切丛的简单例子。我们只讨论形变量子化，把量子化推迟到第 \ref{scbq} 节。

所考虑的例子是式 (\ref{cosu}) 中已引入的复辛流形 $Y$。它由方程\footnote{稍后我们将令 $\j=\hbar j$，最终令 $\hbar=1$，把 $\j$ 换成 $j$。}
\be\label{inco} x^2+y^2+z^2=\j^2 \ee
描述，全纯辛形式为
\be\label{nco}\Omega=\frac{\d x\,\d y}{z}. \ee

该方程与辛形式在 $G=\SO(3,\C)=\PGL(2,\C)$ 对三元组 $x,y,z$ 的显然作用下不变。确实，$Y$ 是 $G$ 的余伴随轨道。它可看作 $G$ 的 Lie 代数 $\g$ 中行列式为 $-\j^2$ 的元素 $T$ 构成的子空间。任何这样的元素可写成
\be\label{pok} T=\begin{pmatrix} z & x+\i y \cr x-\i y & -z\end{pmatrix} \ee
其中 $x,y,z$ 满足式 (\ref{inco})。这样的 $T$ 的特征值为 $\j$ 与 $-\j$，故特征值问题
\be\label{ok}T\begin{pmatrix} q \cr 1\end{pmatrix}=-\j\begin{pmatrix} q \cr 1\end{pmatrix},\qquad T\begin{pmatrix} q' \cr 1\end{pmatrix}=\j\begin{pmatrix} q' \cr 1\end{pmatrix}
\ee
有唯一解，$q,q'\in\C\cup\infty$：
\begin{align}\label{pola} q& = -\frac{x+\i y}{\j+z}=-\frac{\j-z}{x-\i y}\cr
                     q'& = \frac{x+\i y}{\j-z}=\frac{\j+z}{x-\i y}. \end{align}
显然 $q$ 与 $q'$ 在 $\j\leftrightarrow-\j$ 下互换。$q$ 或 $q'$（加上一个无穷远点）参数化 $W=\bCP^1$ 的拷贝，它就是 $T$ 所作用的二维向量空间的射影化。显然 $G$ 将以熟悉的分数线性变换 $q\to(aq+b)/(cq+d)$、$q'\to(aq'+b)/(cq'+d)$（$ad-bc=1$）作用在 $W$ 上。

$Y$ 中没有 $q=q'$ 的点，因为 $T$ 的两个特征向量总不相同。这是唯一的约束；任何一对不同的 $q,q'\in\bCP^1$ 都确定唯一的 $T$，可由式 (\ref{pola}) 解出 $x,y,z$ 得到。$\bCP^1$ 中一点的补集是 $\C$ 的拷贝。故给定 $q$ 时 $q'$ 的可能值参数化一个 $\C$；类似地，给定 $q'$ 时 $q$ 的可能值参数化一个 $\C$。于是 $Y$ 有两个全纯极化：一个的叶由 $q$ 的选择标号（并由满足 $q'\ne q$ 的 $q'$ 参数化），另一个的叶由 $q'$ 标号（并由满足 $q\ne q'$ 的 $q$ 参数化）。把由 $q$ 与 $q'$ 参数化的全纯极化分别记作 $\varPi$ 与 $\varPi'$。

简单计算给出用\footnote{下式中的因子 $1/(q-q')^2$ 并不真正表示 $\Omega$ 的极点，因为如前所述 $q$ 与 $q'$ 总不相同。}$q$ 与 $q'$ 表示辛形式的公式：
\be\label{shortform}\Omega=2\i\j\frac{\d q\d q'}{(q-q')^2}.\ee
也许出乎意料，$q$ 与 $q'$ 的共轭变量可取为相等，因为若
\be\label{momform} p=p'=-\frac{2\i\j}{q-q'}, \ee
则
\be\label{canform}\Omega=\d p\,\d q=\d p'\d q'. \ee

为整体描述此情形，用两个开集 $\U$ 与 $\t\U$ 覆盖 $\bCP^1$，其中 $\U$ 是 $\bCP^1$ 中无穷远点的补集，$\t\U$ 是原点的补集。$\U$ 中的好变量是 $q,q'$；$\t\U$ 中的好变量可由分数线性变换得到：$\t q=-1/q$，$\t q'=-1/q'$。共轭变量为
\be\label{ranform} \t p=\t p'=-\frac{2\i\j}{\t q-\t q'}. \ee
由这些公式可得
\be\label{anform}p=\t q^2\t p+2\i\j\t q=\t q^2\t p-\frac{2\i\j}{q}. \ee
为与一般公式 (\ref{transp}) 比较，注意本例中 $W=\bCP^1$ 的复维为 1，故式 (\ref{transp}) 中的指标 $i,k$ 可省略。我们只用两个开集覆盖 $W$，故 $q,p$ 对应 (\ref{transp}) 中的 $q_\alpha,p_\alpha$，$\t q,\t p$ 对应 $q_\beta,p_\beta$。函数 $\t q^2$ 扮演矩阵 $M$ 的角色；最后，为把 (\ref{anform}) 写成 (\ref{transp}) 的形式，应把 $-2\i\j/q$ 写成 $\i\hbar\partial_q u$，$u$ 为某全纯函数。于是\footnote{函数 $u$ 不是单值的；这是我们用交叠非单连通的开集覆盖 $W$ 所付的代价。第 \ref{shifted} 节的程序会告诉我们用更多、交叠可缩的开集覆盖 $W=\bCP^1$。那时可定义总为单值、但在三重交叠中不满足 $S_{\alpha\beta}S_{\beta\gamma}S_{\gamma\alpha}=1$（除非 $2j\in\Z$）的转移函数。沿这些线索的推导会更长。}$u=-2j\log q$，其中已令 $\j=\hbar j$。然后试图把 $e^u$ 解释为线丛 $\L$ 的转移函数。此时 $e^u=q^{2j}$。它单值当且仅当 $2j\in\Z$，故只有这种情形 $\L$ 才是真正的线丛；对其他 $j$ 值，$\L$ 实际上是线丛的复幂。当 $2j\in\Z$ 时所得的线丛 $\L$ 是熟悉的：在区域 $\t\U$ 中为 1 的截面（在上述 $\t p$ 公式隐含的平凡化下），在区域 $\U$ 中为 $e^u=q^{2j}$（在 $p$ 公式隐含的平凡化下）。该截面在原点有 $2j$ 阶零点、无极点，故 $\L$ 就是通常所说的 $\O(2j)$。无论 $2j\in\Z$ 与否，我们都用记号 $\O(2j)$；要点始终是：作用在 $\O(2j)$ 上的微分算符有定义——我们将写出显式公式——无论 $\O(2j)$ 是否作为线丛存在。

现在讨论空间 $Y$ 上的典范余各向同性 $A$-膜 $\B_\cc$。按一般讨论，$\Hom(\B_\cc,\B_\cc)$ 是作用在 $K_W^{1/2}\otimes\L$ 上的全纯微分算符层。对 $W=\bCP^1$，$K_W\cong\O(-2)$，$K_W^{1/2}\cong\O(-1)$。故 $K_W^{1/2}\otimes\L=\O(2j-1)$。于是预期 $\Hom(\B_\cc,\B_\cc)$ 是 $\bCP^1$ 上作用在 $\O(2j-1)$ 上的全纯微分算符代数。在导出此陈述的某些阶段，我们对 $\Hom(\B_\cc,\B_\cc)$ 作了层化（尽管从物理观点看这一步未必完全自然），但在最终陈述中无需层化。$\Hom(\B_\cc,\B_\cc)$ 就是作用在 $\O(2j-1)$ 上的整体定义的全纯微分算符代数。

然而注意，我们一开始有两个同样好的 $Y$ 的全纯极化，相差 $j\to-j$。若在分析最后一步使用第二个，就会把 $\Hom(\B_\cc,\B_\cc)$ 等同于作用在 $\O(-2j-1)$ 上的整体全纯微分算符代数。因此预期这两个代数同构。确实如此，稍后便知。

$\bCP^1$ 上作用在 $\O(2j-1)$ 上的整体微分算符代数由\footnote{下文令 $\hbar=1$。式 (\ref{throps}) 中的公式在可取这些算符的线性组合并加常数的意义下唯一确定；我们作了方便的选择。}
\begin{align}\label{throps}    t_- & = -\frac{\d}{\d q} \cr
                     t_3 & = q\frac{\d}{\d q}-j+\frac{1}{2} \cr
                      t_+ & = q^2\frac{\d}{\d q}-(2j-1)q \end{align}
生成。这些算符满足 $\mathfrak{sl}(2)$ Lie 代数 $\g$ 的熟悉换位关系：
\be\label{usrel} [t_3,t_\pm]=\pm t_\pm,\qquad [t_+,t_-]=2t_3. \ee
为验证这些算符整体全纯，我们从 $q$ 变换到 $\t q=-1/q$ 并共轭以 $\O(2j-1)$ 的转移函数，检查 $q=\infty$ 处的行为。例如 $t_+=q^2\frac{\d}{\d q}-(2j-1)q$ 变为 $\t q^{2j-1}\left(\frac{\d}{\d\t q}+\frac{2j-1}{\t q}\right)\t q^{-(2j-1)}=\frac{\d}{\d\t q}$，在 $\t q=0$ 处全纯。由于项 $\frac{2j-1}{\t q}$ 的精确系数，极点项抵消了。总之
\begin{align}\label{agrops} t_-& = -\t q^2\frac{\d}{\d\t q}+(2j-1)\t q\cr
                                           t_3& = -\t q\frac{\d}{\d\t q}+j-\frac{1}{2} \cr
                                            t_+ & = \frac{\d}{\d\t q}. \end{align}
算符 $t_+,t_-,t_3$ 代数独立，只差一个关系：把 $\g$ 的二次 Casimir 算符等于一个 $c$-数：
\be\label{eqrel} \frac{1}{2}(t_+t_-+t_-t_+)+t_3^2=j^2-\frac{1}{4}. \ee
它在 $j\to-j$ 下具有预期的对称性。$K_W^{1/2}$ 的平移对得到此对称性至关重要。

算符 $t_+,t_-,t_3$ 是 $Y$ 上全纯函数代数生成元 $x,y,z$ 的量子版本。例如，公式 $p=-2\i j/(q-q')=\i(x-\i y)$ 连同 $p=-\i\d/\d q$ 表明 $t_-=-\d/\d q$ 对应 $x-\i y$。更一般地，梳理上述 $x,y,z$ 与 $p,q$ 的关系，得
\begin{align} \label{correspo} x-\i y &\to t_- \cr
                                             x+\i y &\to t_+ \cr
                                             z& \to t_3. \end{align}
于是 Casimir 关系 (\ref{eqrel}) 是原经典关系 $x^2+y^2+z^2=j^2$ 的量子形变版本。

以上只讨论了形变量子化。第 \ref{scbq} 节将加入拉格朗日 $A$-膜并分析量子化。届时 $Y$ 的一个反全纯对合很重要。若 $\j^2$ 为实，则 $Y$ 有显然的反全纯对合 $\tau:(x,y,z)\to(\bar x,\bar y,\bar z)$。回顾式 (\ref{pola}) 的出发点可知，$\tau$ 交换两个全纯极化 $\varPi$ 与 $\varPi'$（确切地说，$\tau$ 把 $\varPi$ 中给定 $q$ 值的叶映到 $\varPi'$ 中 $q'=-1/\bar q$ 的叶）。同一个复流形 $Y$（仍设 $\j^2$ 实）还有另一个反全纯对合 $\t\tau:(x,y,z)\to(\bar x,-\bar y,-\bar z)$，也把 $\Omega$ 映到 $\bar\Omega$。为写出更清楚的公式，方便的是通过重新定义 $(x,y,z)\to(x,\i y,\i z)$ 把 $\t\tau$ 拉回到 $\tau$。定义 $Y$ 的方程变为
\be\label{nuuxx} x^2-y^2-z^2=j^2, \ee
辛形式仍为 $\Omega=\d x\,\d y/z$。我们考虑这个流形 $Y'$，配以复共轭 $x,y,z$ 的原反全纯对合 $\tau$。当然，纯粹作为复流形，$Y'$ 与 $Y$ 等价，故仍有两个全纯极化。但有个有趣的差别。对 $j^2>0$，两个全纯极化仍被 $\tau$ 交换。但若 $j^2=-s^2$ 为负，则两个全纯极化各自 $\tau$-不变。方程 $x^2-y^2-z^2=-s^2$ 等价于说矩阵
\be\label{uvux}T'=\begin{pmatrix} z & x-y \cr -x-y & -z\end{pmatrix}\ee
的特征值为 $\pm s$。解特征值问题
\be\label{eipro}T'\begin{pmatrix} q \cr 1\end{pmatrix}=-s\begin{pmatrix} q \cr 1\end{pmatrix},\qquad T'\begin{pmatrix} q' \cr 1\end{pmatrix}=s\begin{pmatrix} q' \cr 1\end{pmatrix},\ee
得
\begin{align}\label{horiz} q & = -\frac{x-y}{z+s}=-\frac{z-s}{x+y} \cr
 q'& = -\frac{x-y}{z-s}=-\frac{z+s}{x+y}, \end{align}
它们定义两个全纯极化 $\varPi$ 与 $\varPi'$。由于这些公式完全实，$\varPi$ 与 $\varPi'$ 各自 $\tau$-不变。更确切地说，$\tau$ 通过 $q\to\bar q$ 作用在 $\varPi$ 上，固定实 $q$ 的叶并将其余叶成对交换；对 $\varPi'$ 类似。
